% \begin{table*}[!t]
% \centering
% \resizebox{\textwidth}{!}{%
% \begin{tabular}{@{}lcccc@{}}
% \toprule
% \textbf{Dataset} & \textbf{Arbitrary Personas} & \textbf{Persona-Tailored Questions} & \textbf{Multidimensional} & \textbf{Open-Ended}  \\ \midrule
% RoleLLM \cite{li2023chatharuhi} & \textcolor{purple}{\xmark} & \textcolor{purple}{\xmark} & \textcolor{teal}{\cmark} & \textcolor{teal}{\cmark} \\
% RoleEval \cite{xu2023expertprompting} & \textcolor{purple}{\xmark} & \textcolor{teal}{\cmark} & \textcolor{teal}{\cmark} & \textcolor{purple}{\xmark} \\
% InCharacter \cite{xu2024character} & \textcolor{teal}{\cmark} & \textcolor{purple}{\xmark} & \textcolor{purple}{\xmark}  & \textcolor{teal}{\cmark} \\ \midrule
% \benchmark (Ours) & \textcolor{teal}{\cmark} & \textcolor{teal}{\cmark} & \textcolor{teal}{\cmark}  & \textcolor{teal}{\cmark} \\  \bottomrule
% \end{tabular}%
% }
% \caption{Comparison of existing Persona Agent evaluation benchmarks/frameworks. While several relevant works exists, \benchmark is the only evaluation framework for Persona Agents that enables users to test arbitrary persona agents on their capabilities across multiple dimensions (tasks) through challenging questions tailored to each task and persona.}
% \label{tab:dataset_compare}
% \end{table*}

% \begin{table}[!t]
% \centering
% \resizebox{\columnwidth}{!}{%
% \begin{tabular}{@{}lcccc@{}}
% \toprule
% & \makecell[c]{RoleLLM\\ \cite{li2023chatharuhi}} 
% & \makecell[c]{RoleEval\\ \cite{xu2023expertprompting}}
% & \makecell[c]{InCharacter\\ \cite{xu2024character}} 
% & \makecell[c]{PersonaGym\\ (Ours)} \\ 
% \midrule
% Arbitrary Personas         & \xmark & \xmark & \cmark & \cmark \\
% Persona-Tailored Questions & \xmark & \cmark & \xmark & \cmark \\
% Multidimensional           & \cmark & \cmark & \xmark & \cmark \\
% Open-Ended                 & \cmark & \xmark & \cmark & \cmark \\
% \bottomrule
% \end{tabular}
% }
% \caption{Comparison of existing Persona Agent evaluation benchmarks/frameworks.}
% \label{tab:dataset_compare_transposed_onecolumn}
% \vspace{-15pt}
% \end{table}




\paragraph{Role-Play in LLMs}
Research on LLMs' role-playing capabilities has advanced rapidly. \citet{li2023chatharuhi} enhanced character portrayal through improved prompting and memory extraction, while \citet{xu2024character} examined persona-based decision-making via memory retrieval. \citet{xu2023expertprompting} utilized expert role-playing for QA data generation, and \citet{louie2024roleplay} created a collaborative pipeline where mental health experts provide feedback to guide LLMs in simulating patients. In the counseling domain, \citet{interactive} proposed using dual LLMs to simulate therapist-client interactions. For character development, \citet{characterglm} fine-tuned ChatGLM models for configurable identities, \citet{characterllm} explored profile-based fine-tuning, and \citet{neeko} introduced dynamic LoRA adapters enabling efficient multi-character role-play within a single model.
\vspace{-7.5pt}
\paragraph{Role-Play Evaluation}
Evaluation frameworks for LLM role-playing are emerging. \citet{wang2023rolellm} introduced RoleBench, comprising GPT-generated QA pairs from 100 character profiles. \citet{wang2024incharacter} developed a framework assessing character fidelity through psychological scales and Likert evaluations. \citet{tu2024charactereval} established CharacterEval, a Chinese benchmark containing 1,785 multi-interaction dialogues from novels and scripts. \citet{shen2023roleeval} created RoleEval, a bilingual benchmark with 6,000 multiple-choice questions assessing memorization and reasoning across 300 personas. Table~\ref{tab:dataset_compare_transposed_onecolumn} compares these frameworks with \benchmark, highlighting the necessity of our approach for holistic persona agent evaluation.